% MEASUREMENT_MAP.tex, how each object of the model in MANUSCRIPT.tex is
% observed. Created 10 October 2026 at pass 5 of the plan in TODO.md.
%
% PURPOSE
%   One row per model input and output. Each row says what the object is in
%   the model, what its counterpart is in the Basel standards and the US
%   Call Report, and how it is observed. The FDIC empirics stay in their own
%   document (TODO.md, Decisions).
%
% STRUCTURE THE CHECKER RELIES ON (checks/verify_measurement.py)
%   Rows are written only through
%     \grow{ID}{object}{role in the model}{counterpart and observation}{status}{evidence}
%   IDs are G1, G2, ... in order. Status is one of Rule, Report,
%   Calibration, Unobserved, Output. Evidence is a comma-separated list of
%   quotation or reading IDs of a lit/ record (B3-Q1, LCR-D2, ...), Lean
%   theorems declared under lean/mathlib/, and rows of MANUSCRIPT.tex.
%
% RULES THE CHECKER ENFORCES
%   IDs are well formed, unique and in order. Every status is in the list.
%   Every evidence item resolves. Rule and Report rows cite a quotation from a
%   regulatory document, Calibration rows a quotation from BCVW, Unobserved
%   and Output rows a row of MANUSCRIPT.tex. Every symbol of the coverage
%   list in the checker has a row.

\documentclass[11pt]{article}
\usepackage[T1]{fontenc}
\usepackage[utf8]{inputenc}
\usepackage{mathptmx}
\usepackage{microtype}
\usepackage[a4paper,margin=1.8cm]{geometry}
\usepackage{amsmath,amssymb}
\usepackage{booktabs}
\usepackage{array}
\usepackage{longtable}
\usepackage{url}
\usepackage[hidelinks]{hyperref}
\usepackage{natbib}

\emergencystretch=1.5em

\newcolumntype{P}[1]{>{\raggedright\arraybackslash}p{#1}}

\urlstyle{tt}
\makeatletter
\g@addto@macro\UrlBreaks{\do\.\do\_\do\,\do\-}
\makeatother
\newcommand{\ev}[1]{\path{#1}}

\newcommand{\grow}[6]{#1 & #2 & #3 & #4 & #5 & \ev{#6} \\ \addlinespace}

\title{Measurement Map}
\author{}
\date{Companion to the manuscript skeleton, 10 October 2026}

\begin{document}
\maketitle

\section{Scope}

The map takes each object of the model in \path{MANUSCRIPT.tex} to its
counterpart in the Basel capital framework \citep{bcbs2011}, the Basel
liquidity standard \citep{bcbs2013} and the Call Report instructions for
regulatory capital \citep{ffiec2026}, and says how it is observed. The
regulatory cap is that of \citet{bayraktar2026} (BCVW). Quotation and
reading IDs refer to the \path{lit/} records, where each quotation is checked
verbatim on its page. The rearrangements of the regulatory ratios are proved
in \path{lean/mathlib/MeasurementMap.lean}.

A status of Rule means a regulation sets the value. Report means a
regulatory report records the object for each bank. Calibration means the
model takes a value with no regulatory counterpart. Unobserved means
neither a rule nor a report gives it. Output means the model produces it
and its observation is part of the empirical design.

\section{Map}

\footnotesize
\begin{longtable}{@{}P{0.6cm}P{1.5cm}P{2.7cm}P{5.2cm}P{1.6cm}P{3.6cm}@{}}
\caption{Model objects and their observation.}\label{tab:map}\\
\toprule
ID & Object & Role in the model & Counterpart and observation & Status & Evidence \\
\midrule
\endfirsthead
\toprule
ID & Object & Role in the model & Counterpart and observation & Status & Evidence \\
\midrule
\endhead
\grow{G1}{$x$}{The state, assets over liabilities.}{Total assets over total liabilities. From the reported leverage ratio $\ell$ (Tier 1 over total assets for the leverage ratio), $x=1/(1-\ell)$ holds only if Tier 1 equals book equity and the denominator equals total assets.}{Report}{RCR-Q4, RCR-D3, MeasurementMap.leverage_of_state}
\grow{G2}{$\pi$}{The share of assets in the risky asset.}{Risk-weighted assets over total assets, under risk weights of 0 on the riskless and 1 on the risky asset. With a positive weight on the riskless asset the identification fails.}{Report}{RCR-Q5, RCR-D2, MeasurementMap.rwa_two_assets, MeasurementMap.control_rwa_needs_zero_weight_on_riskless}
\grow{G3}{$X-L$}{Equity, assets less liabilities.}{CET1 capital or Tier 1 capital as reported. Both are net of regulatory adjustments, so they equal book equity only when the adjustments vanish.}{Report}{B3-Q3, RCR-Q1, RCR-Q3, B3-D2}
\grow{G4}{$a_1$}{The floor on equity over risky assets, the solvency branch $(x-1)/a_1$ of the cap.}{The CET1 minimum of 4.5\% of risk-weighted assets, which is BCVW's baseline 0.045, or the Tier 1 minimum of 6.0\%. The bank's own ratio is reported as CET1 or Tier 1 capital over risk-weighted assets. Banks electing the community bank leverage ratio report neither.}{Rule}{BCVW-Q9, BCVW-Q13, B3-Q1, B3-Q2, B3-D1, RCR-Q2, RCR-Q6, RCR-Q7, RCR-D1, MeasurementMap.solvency_iff, MeasurementMap.control_solvency_needs_positive_exposure}
\grow{G5}{$a_2$}{Net 30-day outflows as a share of liabilities, in the liquidity branch $(x-a_2)/a_3$ of the cap.}{A run-off rate on all liabilities. BCVW's baseline 0.05 is the factor for stable insured retail deposits, and less stable retail deposits run off at 10\% or more. The LCR must be at least 100\%. No report read gives the rate.}{Rule}{BCVW-Q10, LCR-Q1, LCR-Q9, LCR-Q10, LCR-D1, MeasurementMap.lcr_iff, MeasurementMap.control_lcr_needs_positive_runoff}
\grow{G6}{$a_3$}{The haircut on the risky asset in high-quality liquid assets, the slope $1/a_3$ of the liquidity branch.}{The riskless asset counts in full, as Level 1 assets do. The standard's haircuts are 15\% (Level 2A) and 25\% or 50\% (Level 2B). BCVW's baseline 0.30 is none of them. The 40\% cap on Level 2 assets and the 75\% cap on inflows are not in the model.}{Calibration}{BCVW-Q11, BCVW-Q12, BCVW-Q13, LCR-Q2, LCR-Q3, LCR-Q4, LCR-Q5, LCR-Q6, LCR-Q7, LCR-Q8, LCR-D2, LCR-D3, MeasurementMap.hqla_haircut, MeasurementMap.cap_iff}
\grow{G7}{$q$}{The requirement on $x-1$, with $z=0$ at $x=1+q$.}{A floor $\ell_0$ on the leverage ratio is the floor $\ell_0/(1-\ell_0)$ on $x-1$, so the 3\% Tier 1 leverage ratio gives $q=3/97$ when Tier 1 equals book equity. This is a reading, and the bundle does not fix $q$. The $z$-volatility does not depend on $q$.}{Rule}{B3-Q7, RCR-Q4, MeasurementMap.leverage_req_iff, MeasurementMap.control_leverage_needs_l0_lt_one, MeasurementMap.zvol2_free_of_q}
\grow{G8}{$\sigma$, $\mu_s$, $r$}{The volatility and drift of the risky asset and the riskless rate.}{No regulatory counterpart. BCVW's baseline is $\sigma=0.08$, $\mu=0.04$, $r=0.02$.}{Calibration}{BCVW-Q14}
\grow{G9}{$\sigma_L$, $c$, $\gamma$, $\mu_L$, $r_L$}{The liability shock, its correlation with the risky return, the net income on liabilities, the liability drift and the deposit rate, with $\mu_L=\gamma+r_L$.}{No regulatory counterpart. BCVW's baseline is $\sigma_L=0.03$, $c=0.20$, $\gamma=0.02$, $\mu_L=0.03$.}{Calibration}{BCVW-Q14, M9}
\grow{G10}{$\rho_L$}{The liability-net discount rate of the reduced problem.}{Formed from the shareholder discount rate, BCVW's baseline $\rho=0.12$. No regulatory counterpart.}{Calibration}{BCVW-Q14}
\grow{G11}{$\kappa$}{The proportional issuance cost.}{No regulatory counterpart. BCVW's baseline is $\kappa=0.01$ for dilution and $\kappa'=0.02$ per unit issued.}{Calibration}{BCVW-Q14}
\grow{G12}{$K$}{The fixed issuance cost, charged per unit of liabilities.}{Not observed. A cost fixed in dollars would break the reduction to $x$. Locally identified together with $\lambda_S$ from the barrier and the trigger.}{Unobserved}{BCVW-Q4, K2, M5, M6}
\grow{G13}{$\lambda_S$}{The asymmetry parameter of the shortfall penalty.}{Not observed. The variance ratio carries no information about it. Locally identified together with $K$ from the barrier and the trigger.}{Unobserved}{K1, K2, M3, M6, M11}
\grow{G14}{$y^*$}{The dividend barrier.}{The state at which the bank pays dividends. Under the conservation buffer a bank with CET1 between 4.5\% and 7.0\% of risk-weighted assets has its distributions restricted, so observed payouts reflect the rule as well as the barrier.}{Output}{M4, B3-Q4, B3-Q5, B3-Q6, B3-D3}
\grow{G15}{$x_L$, $y_{\text{post}}$}{The recapitalisation trigger and the injection target.}{The state just before and just after an equity issue.}{Output}{M5, M6, M7, M10}
\grow{G16}{$\zeta$, $\nu_1^2$, $\nu_3^2$}{The $z$-volatility and its limits at the distress boundary and deep in surplus.}{The variance of changes in $z=\log((x-1)/q)$ near the distress boundary and deep in surplus. It does not depend on $q$.}{Output}{M2, M3, M8, MeasurementMap.zvol2_free_of_q}
\grow{G17}{$\nu_1^2/\nu_3^2$}{The deficit-to-surplus variance ratio, read in discrete time as $\lambda_V^4$.}{The ratio of the two variances of G16. Reading it as $\lambda_V^4$ is a definition.}{Output}{K1, M1, M3}

\bottomrule
\end{longtable}
\normalsize

\section{The ratio constraints in the state}

With equity $X-L$, risk-weighted assets $\pi X$, high-quality liquid assets
$X-a_3\pi X$ and net outflows $a_2L$, the two ratio constraints of BCVW are
the two branches of the cap in $x=X/L$. For $a_1,a_2,a_3,\pi,X,L>0$,
\[
a_1\le\frac{X-L}{\pi X}\iff\pi x\le\frac{x-1}{a_1},\qquad
1\le\frac{X-a_3\pi X}{a_2L}\iff\pi x\le\frac{x-a_2}{a_3},
\]
and both hold if and only if $\pi x\le u(x)=\min\{(x-1)/a_1,(x-a_2)/a_3\}$
(\path{MeasurementMap.solvency_iff}, \path{MeasurementMap.lcr_iff},
\path{MeasurementMap.cap_iff}). At $\pi=0$ the left side of the first
equivalence is undefined and the equivalence fails, and at $a_2=0$ the same
holds for the second (the two controls). For $\ell_0<1$ and $X,L>0$,
$\ell_0\le(X-L)/X\iff\ell_0/(1-\ell_0)\le x-1$
(\path{MeasurementMap.leverage_req_iff}).

\section{Pending}

\begin{itemize}
\item The FDIC field names used by the empirical scripts (\path{RWAJT},
  \path{RBCT1J}, \path{RBCRWAJ}, \path{ASSET}, \path{LIAB}, \path{EQCDIV},
  \path{EQCSTKRX}) do not appear in the Call Report instructions read. Their
  link to the items of G1 to G4, G14 and G15 needs the FDIC data dictionary.
\item Rows G1 to G3 assume regulatory capital equals book equity. Items
  12 to 18 and 27 to 29 of Schedule RC-R, which give the adjustments, were
  not read.
\item The tier that $a_1$ stands for is not fixed (B3-D1).
\end{itemize}

\bibliographystyle{plainnat}
\bibliography{refs}

\end{document}
